\scanpage{331}
\mspage{70}
% archive sheet number: 330
% ms. 1–2
Donc il y a deux vérifications à faire (où je suppose à nouveau $\Sigma$ quelconque
% ms. 3
1°) Soit $u\in\Sigma$, \inserted{$u:a\longrightarrow b$,} supposons $\Sigma$ inversible \inserted{dans $X$}. Alors
% ms. 4
\[\mathrm{ch}(u)^{\unsure{\circ|o}}=\mathrm{ch}(u^{-1})\quad\text{dans }\Pi_\Sigma(X)\]
% ms. 5
i.e. les deux chemins de $b$ vers
% ms. 6–7
\diagram{\begin{tikzcd}[ampersand replacement=\&,row sep=small]
b \& a\arrow[l,very thick,"u"']\arrow[l,phantom,"{\deleted{\unsure{\to}}}"{description,very near start}] \& \text{de type }\Delta_-\\
b\arrow[r,"u^{-1}"] \& a \& \text{de type }\Delta_+
\end{tikzcd}}
% ms. 8–10
sont homotopes dans $\underline{\mathrm{Ch}}_\Sigma(b,a)$. Or on sait déjà, \deleted{que} comme $\unsure{k|K}:X\longrightarrow\Pi_\Sigma(X)$ est un foncteur, que
% ms. 11
\[\deleted{\unsure{h|\mathrm{hot}}(\mathrm{ch}(u^{-1}))=(\unsure{h|\mathrm{hot}}(\mathrm{ch}(u))^{-1}\quad\text{dans }\Pi_\Sigma(X)}\]
% ms. 12
\deleted{donc il faut prouver que $\mathrm{ch}(u)^{-1}={}$}
% ms. 13
\[\mathrm{hot}\ \mathrm{ch}(u^{-1})=(\mathrm{hot}\ \mathrm{ch}(u))^{-1}\quad\text{dans }\Pi_\Sigma(X),\]
% ms. 14–18
où $\mathrm{hot}(\bullet)$ désigne la classe d'homotopie d'une \unsure{chaîne} (ici \unsure{chem.|chemin|dans} dans $\underline{\mathrm{Ch}}_\Sigma(b,a)$ \deleted{\unsure{$f$}} et dans $\underline{\mathrm{Ch}}_\Sigma(a,b)$ respectivement. Donc l'énoncé à prouver c'est que $\mathrm{hot}(\mathrm{ch}(u)^\circ)$ est égal à l'\deleted{\unsure{id.|identité|idem}}\inserted{\uline{inverse}} de $\mathrm{hot}(\mathrm{ch}(u))$ dans $\Pi_\Sigma$. C'est donc ce qu'on a prouvé pour \uline{\unsure{tt|$\forall$}} $u\in\mathrm{Fl}(X)$, dans le cas de $\Pi(X)$. L'argument doit marcher, du moment qu'il a un sens \unsure{i.e.|ici,} que $\mathrm{ch}(u)^\circ$ est bel et bien un $\Sigma$-chemin, i.e. $u\in\Sigma$.

% ms. 24–25
2°) La formule (2) signifie que pour tt chemin $c$ de type $\Delta_+^n$ (\unsure{donc} tautologiquement un
